Reinforced concrete (RC) structures are widely used in critical infrastructure systems; however, their long-term durability is strongly affected by degradation mechanisms such as carbonation-induced corrosion of steel reinforcement. Carbonation reduces the alkalinity of concrete, enabling corrosion initiation and progressive loss of structural capacity, which may compromise safety, serviceability, and lifecycle performance \citep{wang2024review, aslaniprobabilistic2020}. As infrastructure networks age and maintenance resources become increasingly constrained, the ability to accurately forecast the remaining useful life (RUL) of RC structures has become a central challenge in structural engineering and risk-informed asset management \citep{ellingwood2016introduction}.

Traditional service life prediction approaches rely on physics-based carbonation depth models and deterministic limit-state formulations. While these models provide valuable mechanistic insight, they often struggle to capture the combined effects of material variability, environmental uncertainty, stochastic deterioration, and epistemic model error \citep{pinto2022service}. Recent advances in probabilistic durability modeling have therefore emphasized the integration of uncertainty quantification, stochastic deterioration processes, and data-driven calibration to improve predictive reliability \citep{wang2024review}.

In parallel, machine learning (ML) and surrogate modeling techniques have gained increasing attention for predicting carbonation depth, structural degradation, and lifecycle performance of RC systems. Comparative studies have demonstrated the ability of artificial neural networks and other ML models to outperform conventional regression-based carbonation models when trained on experimental and field data \citep{couto2025machine}. Nevertheless, purely data-driven models often exhibit limited extrapolation capability, reduced interpretability, and sensitivity to sparse training data, which restricts their direct application in safety-critical reliability assessments.

To address these limitations, surrogate modeling frameworks grounded in uncertainty quantification have emerged as promising alternatives. Polynomial Chaos Expansion (PCE) provides an efficient spectral representation of stochastic system responses, enabling rapid reliability evaluation while preserving probabilistic interpretability \citep{lim2021distribution, lamouri2024reliability}. PCE-based metamodels have demonstrated strong performance in structural reliability analysis by significantly reducing computational cost relative to Monte Carlo simulation, while maintaining high accuracy across multiple failure domains.

However, global surrogate models such as PCE may exhibit reduced fidelity in highly nonlinear or locally irregular regions of the response surface. To overcome this limitation, local Gaussian-process–based approximation strategies have been proposed to refine predictions in critical regions of the stochastic domain. Local surrogate refinement techniques have shown improved robustness and accuracy in structural reliability problems, particularly when combined with global spectral surrogates \citep{search17surrogate2017}. 

In this context, the present study proposes a hierarchical modeling framework that integrates Polynomial Chaos Expansion with a Gaussian Local Approximation Model (GLAM). The rationale for this sequential strategy is twofold. First, PCE is employed as a global emulator to capture the dominant stochastic structure of resistance and demand fields while ensuring computational efficiency. Subsequently, GLAM is applied as a localized correction layer to enhance predictive accuracy in regions where failure probability sensitivity is highest. This ordered coupling leverages the complementary strengths of global spectral approximation and local probabilistic refinement, yielding a balanced trade-off between efficiency, robustness, and accuracy.

Building on this framework, the proposed methodology is applied to the forecasting of the remaining useful life of reinforced concrete structural elements subjected to carbonation-induced corrosion. The approach integrates stochastic resistance modeling, time-dependent carbonation progression, probabilistic reliability assessment, and emulator-based uncertainty propagation. To verify the proposed methodological pipeline, an adapted version of a benchmark reliability problem is considered, following the general validation strategy proposed by \citet{pires_reliability_2024}. 

The main contributions of this work are threefold: (i) the development of a hybrid global–local surrogate modeling architecture for time-dependent structural reliability, (ii) the integration of carbonation-driven deterioration into a probabilistic RUL forecasting framework, and (iii) the demonstration of improved predictive efficiency and accuracy relative to conventional Monte Carlo and single-surrogate approaches. The results aim to advance reliability-based lifecycle assessment of RC infrastructure and support risk-informed maintenance and decision-making.






